\section{Introduction}
\label{sec:intro}

The advent of large-scale pre-training has fundamentally altered the paradigm of deep learning. Instead of training model architectures from scratch, practitioners commonly adapt a single pre-trained base model to multiple distinct downstream tasks, resulting in a set of task-specific expert models. While multi-task ensembling provides an effective route for model deployment, storing and serving multiple large neural networks concurrently introduces prohibitive computational and storage overhead. 

To address this challenge, post-hoc model merging has emerged as a promising alternative, seeking to blend the weights of multiple fine-tuned models directly in parameter space without requiring joint retraining \cite{wortsman2022model, ilharco2022editing}. Early heuristics, such as Task Arithmetic \cite{ilharco2022editing} and TIES-Merging \cite{yadav2023ties}, employ static, coordinate-wise combinations. Although simple, these static approaches fail to capture the complex, layer-specific functional properties of deep neural networks, often leading to severe representation interference and performance degradation under task conflicts.

To overcome the limitations of static merging, adaptive ensembling paradigms have been proposed to dynamically tune merging coefficients. For example, Online AdaMerging \cite{yang2023adamerging} adapts merging coefficients at test-time by minimizing the prediction entropy of local test streams. However, as demonstrated by recent audits \cite{suitemerge2025}, unsupervised test-time adaptation (TTA) over small local streams is highly vulnerable to \emph{transductive overfitting} and trivial class collapse, where the model minimizes entropy by predicting a single class. As an alternative, Offline Few-Shot Validation Tuning (OFS-Tune) leverages a tiny, labeled calibration dataset (e.g., $M=10$ samples per task) to optimize merging parameters. Yet, this approach introduces its own failure mode: when optimizing an independent continuous parameter for every layer and task (a search space of size $K \times L$), the extremely high capacity of the hypothesis class relative to the minuscule calibration set leads to severe generalization gaps and high-frequency parameter oscillations.

In this work, we approach this dilemma through the lens of \textbf{Statistical Learning Theory}. We present \textbf{Rademacher-Bounded Polynomial Merging (RBPM)}, a mathematically rigorous framework that provides the first formal generalization guarantees for adaptive model merging. We treat the layer-wise merging coefficients across network depth as a continuous global trajectory, rather than a set of decoupled variables. By projecting this high-dimensional trajectory space onto a low-degree polynomial subspace of degree $d \ll L$, we demonstrate that we can filter out high-frequency layer-specific noise.

Using the tools of statistical learning theory, we derive tight **Rademacher Complexity Bounds** for this polynomial trajectory class. We prove that constraining the coefficients to a smooth polynomial trajectory scales the empirical Rademacher complexity with the polynomial degree $d$ rather than the network depth $L$, shrinking the complexity by a factor of $\mathcal{O}(\sqrt{L / \log(d)})$. This theoretical guarantee directly bounds the generalization gap between calibration and test performance. Furthermore, we incorporate an analytical Rademacher regularization penalty (corresponding to an $L_1$ constraint on the polynomial coefficients) directly into the few-shot optimization objective to strictly control the capacity of the merged model.

Our empirical evaluations on a 12-layer deep CNN architecture across four distinct visual classification tasks confirm our theoretical insights:
\begin{enumerate}
    \item \textbf{Robustness Against Overfitting:} Our proposed RBPM framework ($\lambda_{\text{rad}} = 0.01$) achieves a robust average test accuracy of 38.85\%, vastly outperforming the Static Uniform baseline (29.05\%) and the Offline Unconstrained Few-Shot Tuning baseline (32.75\%). This confirms that smooth global trajectories combined with statistical capacity control filter out local validation noise and prevent overfitting.
    \item \textbf{Efficacy of Trajectory-Constrained TTA:} We show that unsupervised online adaptation methods (such as Online AdaMerging and Online PolyMerge) perform well when using robust batch statistics on test streams. Interestingly, Online PolyMerge ($d=2$) achieves higher performance than the unconstrained Online AdaMerging, demonstrating that polynomial trajectory constraints act as a powerful regularizer not only offline but also online under entropy minimization. However, our supervised offline few-shot calibrated RBPM still achieves the overall highest test accuracy by leveraging the few-shot labels to resolve parameter conflicts.
    \item \textbf{Generalization Gap Control:} By sweeping the Rademacher regularization coefficient $\lambda_{\text{rad}}$, we demonstrate precise control over the generalization gap, reducing the gap to a tight --1.35\% at the optimal peak ($\lambda_{\text{rad}}=0.01$), and recovering the Uniform baseline exactly under strong regularization ($\lambda_{\text{rad}}=1.0$), establishing a provable bridge between learning theory and parameter-space ensembling.
\end{enumerate}
Our work brings statistical rigor to a field dominated by heuristics, demonstrating that theoretical guarantees can directly guide the design of robust, high-performance model merging algorithms.